% Ce qui se passe à l'expiration d'une clé populaire, selon la stratégie (schéma).
\documentclass[border=6pt]{standalone}
\usepackage{sd}
\begin{document}
\begin{tikzpicture}[x=1mm, y=1mm, every node/.append style={font=\sffamily\small}, note/.append style={font=\sffamily\small}]
  \foreach \nom/\y in {{naïve}/0, {verrou}/-20, {stale-while-revalidate}/-40} {
    \node[axe, anchor=east] at (-3,\y) {\nom};
  }
  \draw[Dink, line width=1.2pt] (40,6) -- (40,-46);
  \node[note, anchor=south] at (40,6) {la clé expire};
  % naïve : beaucoup de recalculs en même temps
  \foreach \i in {0,...,9} \fill[DredBg, draw=Dred, line width=0.5pt] (40+\i*0.9,{-3+\i*0.6}) rectangle (60+\i*0.9,{-2+\i*0.6});
  \node[note, text=Dred, anchor=west, align=left] at (72,1) {chaque lecteur recalcule :\\80 recalculs par expiration (mesuré)};
  % verrou : un recalcul, les autres attendent
  \fill[DgreenBg, draw=Dgreen, line width=0.8pt] (40,-22) rectangle (60,-18);
  \node[note, text=Dgreen] at (50,-20) {1 recalcul};
  \fill[DamberBg, draw=Damber, line width=0.5pt] (40,-27) rectangle (60,-24);
  \node[note, text=Damber, anchor=west, align=left] at (62,-25.5) {les autres lecteurs attendent :\\99e centile 87 à 89 ms (mesuré)};
  % SWR : périmée servie, un rafraîchissement en arrière-plan
  \fill[DslateBg, draw=Dslate, line width=0.5pt] (40,-44) rectangle (60,-41);
  \node[note, anchor=west] at (61,-42.5) {valeur périmée servie à tous : 99e centile 5 à 6 ms (mesuré)};
  \fill[DgreenBg, draw=Dgreen, line width=0.8pt] (40,-38.5) rectangle (60,-35);
  \node[note, text=Dgreen] at (50,-36.7) {en fond};
  \node[note, text=Dgreen, anchor=west] at (61,-36.7) {1 recalcul en arrière-plan};
  \draw[ligne] (0,-50) -- (120,-50);
  \foreach \t/\l in {40/{0 ms}, 60/{100 ms}} {\draw[ligne] (\t,-50) -- (\t,-51.5); \node[note, below] at (\t,-51.5) {\l};}
  \node[axe, anchor=north] at (60,-56) {temps (le recalcul dure 100 ms)};
\end{tikzpicture}
\end{document}
